% Part: computability
% Chapter: recursive-functions
% Section: general-recursive-functions

\documentclass[../../../include/open-logic-section]{subfiles}

\begin{document}

\olfileid{cmp}{rec}{gen}
\olsection{Fungsi Rekursif Umum}

Ana cara liya kanggo ngentukake himpunan fungsi total.
Fungsi total $f(x,\vec z)$ diarani \emph{reguler}
yen kanggo saben runtutan wilangan asli $\vec z$,
ana $x$ kang ndadekake $f(x,\vec z) = 0$.
Kanthi tembung liya, fungsi reguler yaiku persis
fungsi kang bisa dikenani panggolekan tanpa wates
lan asile tetep fungsi total. Tanpa ngowahi
himpunan fungsi total kang diasilake, kita bisa
matesi panggolekan tanpa wates mung kanggo
fungsi reguler:

\begin{defn}
\ollabel{defn:general-recursive}
Himpunan \emph{fungsi rekursif umum} yaiku
himpunan paling cilik saka fungsi ing wilangan asli
menyang wilangan asli (kanthi aritas warna-warna)
kang ngemot fungsi nol, penerus, lan proyeksi,
lan katutup tumrap komposisi, rekursi primitif,
lan panggolekan tanpa wates kang ditrapake
marang fungsi \emph{reguler}.
\end{defn}

Cetha yen saben fungsi rekursif umum iku total.
Bedane \olref{defn:general-recursive} lan
\olref[par]{defn:recursive-fn} yaiku ing definisi
kang pungkasan, fungsi rekursif parsial oleh
digunakake ing langkah-langkah antarane; siji-sijine
syarat yaiku fungsi kang diasilake ing pungkasan
kudu total. Mula tembung 'umum', tinggalan
sajarah, iku sebutan kang kurang trep; yen
katon saka definisine,
\olref{defn:general-recursive} malah \emph{kurang}
umum tinimbang \olref[par]{defn:recursive-fn}.
Nanging untunge, bedane iki mung katon wae;
sanajan definisine beda, himpunan fungsi
rekursif umum lan himpunan fungsi rekursif
iku padha persis.

\end{document}
